% Section VII -- Results. Budget: 2.4 col incl. Fig. 3, Fig. 4, Table II (PAPER_PLAN.md).
% Every number: generated tables in ../generated/ (pipeline, session medians per D26) or PAPER_PLAN final numbers.
\section{Results}\label{sec:results}

\begin{figure}[t]
\centering
\includegraphics[width=\columnwidth]{fig_a3_cycles.pdf}
\caption{Per-layer cycles: model, RTL simulation and measured on the KV260
for both networks. All three agree exactly on every layer.}
\label{fig:cycles}
\end{figure}

\textbf{Three-way agreement.} Every layer's cycle count is identical in the
model, in RTL simulation and on the KV260 (Fig.~\ref{fig:cycles}): 16,436 and
104,247 cycles per LeNet-5 and CIFAR-10 image. On the board, each layer took
a single cycle count over all 10,000 test images of each network (measured
on the KV260, one session), as in RTL simulation.
% src: tab_a3_cycles (model = RTL = KV260 per layer, 0.00 % error); tab_layer_spread (10,000 images, 1 distinct value, spread 0)
The 300 random-shape jobs give the same result on the board as in RTL
simulation: all 286 legal jobs are bit- and cycle-exact and all 14 illegal
jobs are refused (Fig.~\ref{fig:shapes}). Cycle counts are also identical at
all 7 clocks of the clock sweep. % src: tab_shapes (KV260 row); PAPER_PLAN final numbers ("Cycles across 7 clocks: identical at every clock", hw_b2_cycles.csv)

\begin{figure}[t]
\centering
\includegraphics[width=\columnwidth]{fig_shapes_cycles.pdf}
\caption{Random-shape jobs: cycles predicted by the model vs.\ measured in
RTL simulation and on the KV260 (every point lies on $y = x$).}
\label{fig:shapes}
\end{figure}

\textbf{Correctness.} The board's logits differ from the golden model on
0 of 10,000 images per network, so its accuracy equals the INT8 reference:
98.79\,\% (LeNet-5) and 78.52\,\% (CIFAR-10), against 98.78\,\% and
78.56\,\% in FP32 (model). A 30-minute soak ran 3.65M inferences with no
logit, cycle or status error and one cycle count per network, at a die
temperature of at most 37.7\,$^\circ$C (measured on the KV260).
% src: tab_a1_accuracy; tab_soak (3,651,181 jobs, 0 mismatches, 1 distinct TOTAL_CYC, 37.7 C)

\begingroup
\renewenvironment{table}[1][t]{\begin{table*}[#1]}{\end{table*}}
\input{../generated/tab_a5_cpu}
\endgroup

\textbf{Latency and throughput.} Accelerator compute takes 65.7\,$\mu$s and
417.0\,$\mu$s per image (measured on the KV260), with no spread across
images or sessions (Table~\ref{tab:cpu}); this is 15,210 and 2,398 images/s,
or 12.7 and 21.6 GOPS, 39.6\,\% and 67.3\,\% of the array's 32 GOPS peak,
matching the utilization of Sec.~\ref{sec:model}.
% src: tab_a5_cpu (65.7 [65.7, 65.7]; 417.0 [417.0, 417.0]); tab_efficiency (15,210 / 2,398 img/s; 12.67 / 21.55 GOPS; 39.6 / 67.3 % of 32.0 peak)
End to end, the Python host adds most of the time: the median is
224.5\,$\mu$s and 578.5\,$\mu$s, with p99 at 239.1\,$\mu$s and 593.2\,$\mu$s.
Determinism therefore holds for accelerator compute; the end-to-end spread
comes from the host software.
% src: tab_percentiles (fast host path p50 / p99)

\textbf{Comparison.} The best CPU baseline is ONNX Runtime INT8 with
4 threads for both networks: 696.9\,$\mu$s and 1,222.9\,$\mu$s end to end, so
the accelerator is 3.10$\times$ and 2.11$\times$ faster (Table~\ref{tab:cpu}).
The CPU's end-to-end time also includes input preprocessing, which ours does
not; on compute alone the ratio is 9.41$\times$ and 2.58$\times$.
% src: tab_a5_cpu (ORT INT8 x4 e2e 696.9 / 1,222.9; speedups e2e 3.10x / 2.11x, comp 9.41x / 2.58x; note on preprocessing)
The DPU takes 345.6\,$\mu$s and 352.8\,$\mu$s end to end (one session): the
accelerator is 1.5$\times$ faster on LeNet-5, and the DPU is 1.6$\times$
faster on CIFAR-10, where its 2,048 MAC lanes at 300\,MHz outweigh its
fixed overheads (our array has 64 lanes at 250\,MHz).
% src: tab_a5_cpu (DPU e2e p50 345.6 / 352.8; DPU/accel e2e 1.54x / 0.61x); tab_efficiency (MAC lanes 64 / 2,048; lane clock 250 / 300); PAPER_PLAN ("1.6x slower")
% UNSURE: "outweigh its fixed overheads" is an interpretation, not a measured decomposition.

\input{../generated/tab_power_compare}

\textbf{Power and energy.} Table~\ref{tab:powercmp} reports SOM-rail power
(INA260). With no overlay the SOM idles at 3.207\,W; loading our design
raises idle to 3.44\,W, the DPU overlay to 4.36\,W. Running adds
0.251\,W and 0.263\,W for the accelerator, against 0.355\,W and 0.455\,W for
the DPU. Per image, the accelerator uses 0.065\,mJ and 0.161\,mJ, the DPU
0.124\,mJ and 0.166\,mJ, and the best CPU baseline 0.454\,mJ and 0.895\,mJ:
about half the DPU's energy on LeNet-5, the same on CIFAR-10, and 5.5--7$\times$
less than the CPU.
% src: tab_power_compare (3.207; 3.442 / 3.433; 4.362 / 4.361; dP 0.251 / 0.263 vs 0.355 / 0.455; E 0.065 / 0.161 vs 0.124 / 0.166); tab_b1_power (CPU 0.454 / 0.895 mJ); PAPER_PLAN ("5.5-7x less energy")
Most of the accelerator's $\Delta P$ is the host loop: the control phase,
which runs the same loop without accelerator compute, already draws
0.222\,W and 0.211\,W, leaving 0.029\,W and 0.051\,W for the PL work.
% src: tab_b1_power (dP control 0.222 / 0.211; dP accel - control 0.029 / 0.051)

\textbf{Repeatability and clock sweep.} Across the three sessions the
end-to-end medians stay within [223.3, 225.4]\,$\mu$s and [576.9,
580.5]\,$\mu$s, and the energy per image varies by a CV of 3.0\,\% (LeNet-5)
and 0.4\,\% (CIFAR-10). % src: tab_repeatability (B3 e2e fast ranges); tab_b1_power (E_sys CV 3.0 / 0.4 %)
Over 7 clocks from 100 to 250\,MHz, latency scales as $1/f$ with constant
cycles, and SOM-rail power fits $P = 3.424\,\mathrm{W} + 1.095\,\mathrm{mW/MHz}
\cdot f$ ($R^2 = 0.9964$); a second sweep gave a slope 8\,\% different, so
we treat the split between static and dynamic power as approximate.
% src: tab_b2_fit (P_accel absolute: 3.424, 1.095, R2 0.9964, n 7; k differs 8 % from the ascending sweep); EXPERIMENTS.md B2 (100-250 MHz)
